\section{References}
\label{app:references}

\subsection{Sacred Scripture and Ancient Texts}

\begin{itemize}
\item \textbf{Clementine Vulgate.} \emph{Biblia Sacra Vulgatae Editionis} (Sixto-Clementine text, 1592; eBible \texttt{latvuc}; public domain). Matt 27--28; Mark 15--16; Luke 4, 19, 23--24; John 11, 19--20; Acts 10--11, 19; Rom 6.
\item \textbf{Douay--Rheims Bible.} Bishop Richard Challoner revision (Rheims NT 1582; Douay OT 1609--1610; Challoner 1749--1752; public domain). Controlling English Scripture text.
\item \textbf{King James Version (1611), Revised Version (1895), and World English Bible (Catholic Edition).} Public-domain English translations cited for lexical comparison at Luke 19:20, 24:12, John 11:44, 20:5--9, and Acts 19:12.
\item \textbf{Robinson, Maurice A., and William G. Pierpont, eds.} \emph{The New Testament in the Original Greek: Byzantine Textform} (RP2018 Unicode CSV with inline Nestle--Aland divergence notes; public domain).
\item \textbf{Holmes, Michael W., ed.} \emph{The Greek New Testament: SBL Edition} (SBLGNT v1.2, Faithlife snapshot). Luke 19:20, 23:53, 24:10--13.
\item \textbf{Tauber, James, ed.} \emph{MorphGNT SBLGNT} (v6.12 lemma-field dataset; CC BY-SA 3.0). Complete New Testament lemma counts.
\end{itemize}

\subsection{Lexicography and Textual Criticism}

\begin{itemize}
\item \textbf{Thayer, Joseph Henry.} \emph{A Greek-English Lexicon of the New Testament}, corrected edition (New York: American Book Company, undated title page; 1889; public domain). Entries s.v. \latin{entyliss\={o}} (p. 219), \latin{othonion} (p. 439), and \latin{soudarion} (p. 581).
\item \textbf{Westcott, Brooke Foss, and Fenton John Anthony Hort.} \emph{The New Testament in the Original Greek: Introduction and Appendix} (Cambridge and London: Macmillan, 1882; public domain). Appendix, \emph{Notes on Select Readings}, p. 71, on Luke 24:12.
\end{itemize}

\subsection{Patristic, Medieval, and Catholic Commentators}

\begin{itemize}
\item \textbf{Ambrose of Milan, St.} \emph{Expositio evangelii secundum Lucam}, Book X, \S 179. Migne, \emph{Patrologia Latina} 15 (via \emph{Corpus Corporum}; public-domain base text).
\item \textbf{Aquinas, St. Thomas.} \emph{Super Evangelium S. Ioannis Lectura}, cap. 20, lect. 1. Latin text edited by Roberto Busa and Enrique Alarc\'on (\emph{Corpus Thomisticum}, based on Turin 1952 printing).
\item \textbf{Aquinas, St. Thomas.} \emph{Catena Aurea: Commentary on the Four Gospels, Collected out of the Works of the Fathers}, Vol. IV: \emph{St. John}, new edition (Oxford: James Parker and Co., 1874; historical English translation; public domain). Cited on John 11:44 (pp. 382--383) and John 20:1--10 (pp. 593--596).
\item \textbf{Augustine of Hippo, St.} \emph{In Iohannis Evangelium Tractatus}, Tractates 49.24, 120.7--9, and 121.1. English translation by John Gibb and James Innes in Philip Schaff, ed., \emph{Nicene and Post-Nicene Fathers}, First Series, Vol. VII (New York: The Christian Literature Company, 1888; public domain); Latin text of Tractate 120.9 checked via \texttt{augustinus.it}.
\item \textbf{Augustine of Hippo, St.} \emph{De Consensu Evangelistarum} (\emph{The Harmony of the Gospels}), Book III, chs. 23--25 (\S\S 59, 68--70). English translation by S. D. F. Salmond, edited by M. B. Riddle, in \emph{Nicene and Post-Nicene Fathers}, First Series, Vol. VI (public domain).
\item \textbf{Bede the Venerable, St.} \emph{In Lucae evangelium expositio}, Book VI, on Luke 24:12. Migne, \emph{Patrologia Latina} 92, col. 625B (via \emph{Corpus Corporum}; public-domain base text).
\item \textbf{Cornelius a Lapide, S.J.} \emph{The Great Commentary of Cornelius \`a Lapide: The Gospels}, Vols. 5 (John I--XI) and 6 (John XII--XXI), translated by Thomas W. Mossman et al. (Edinburgh: John Grant, 1908; public domain). Cited on John 11:44 (vol. 5, pp. 409--410) and John 20:1--11 (vol. 6, pp. 255--257).
\item \textbf{Cyril of Alexandria, St.} \emph{Commentary on the Gospel according to S. John}, Vol. II, Book XII (on John 20:1--9), translated by Thomas Randell, \emph{A Library of Fathers of the Holy Catholic Church} (London: Walter Smith, 1885; public domain), pp. 649--651.
\item \textbf{Gregory the Great, St.} \emph{Homiliae in Evangelia}, Homily 22 (\S\S 3--5). Migne, \emph{Patrologia Latina} 76 (via \emph{Corpus Corporum}; public-domain base text); additional excerpt on John 20:1--9 via Aquinas, \emph{Catena Aurea in Ioannem} (pp. 595--596).
\item \textbf{John Chrysostom, St.} \emph{Homilies on the Gospel of St. John}, Homilies 85.4, 86.1, and 87. English translation by Charles Marriott, edited by Philip Schaff, \emph{Nicene and Post-Nicene Fathers}, First Series, Vol. XIV (New York: Christian Literature Co., 1889; public domain).
\item \textbf{Theophylact of Ohrid, Bl.} \emph{Enarratio in Evangelium Ioannis}, on John 20. Migne, \emph{Patrologia Graeca} 124, facing Latin column (Paris, 1864; public domain), cols. 289--291; moral typology excerpted via Aquinas, \emph{Catena Aurea in Ioannem} (p. 596).
\end{itemize}

\subsection{Magisterial and Liturgical Witnesses}

\begin{itemize}
\item \textbf{Catechism of the Council of Trent (Roman Catechism).} Translated into English by the Rev. J. Donovan (Maynooth, 1829; Baltimore: Lucas Brothers; public domain). Part I, Article III (p. 41) and Article V (pp. 53--54).
\item \textbf{Missale Romanum (1962).} Typical edition of 1962 (corroborated against the public-domain Pustet Ratisbon 1862 printing). Cited for \emph{Victimae paschali laudes} (Easter Sunday and Octave) and the Gospel appointment of John 20:1--9 for Easter Saturday (\latin{Sabbato in albis}).
\item \textbf{Ordo Lectionum Missae (1981).} Second typical edition, no. 42. Cited for the Easter Sunday principal Gospel appointment of John 20:1--9.
\item \textbf{Caswall, Edward.} \emph{Lyra Catholica: Containing All the Breviary and Missal Hymns} (London: James Burns, 1849; public domain), pp. 233--234. Literary (non-liturgical) English translation of \emph{Victimae paschali laudes}.
\end{itemize}

\subsection{Rabbinic, Classical, and Early-Modern Historical Witnesses}

\begin{itemize}
\item \textbf{Babylonian Talmud, Tractate \emph{Berakhot}.} Folios 42a and 52b. Translated by Abraham Cohen, \emph{The Babylonian Talmud: Tractate Berakot} (Cambridge: Cambridge University Press, 1921; public-domain text via Sefaria).
\item \textbf{Lightfoot, John.} \emph{Horae Hebraicae et Talmudicae: Hebrew and Talmudical Exercitations upon the Gospels, the Acts, some Chapters of St. Paul's Epistle to the Romans, and the First Epistle to the Corinthians}, edited by Robert Gandell, 4 vols. (Oxford: Oxford University Press, 1859; public domain). Vol. 3 on John 11:44 and John 20 (absent lemma on John 20:7).
\item \textbf{Mishnah, Tractate \emph{Berakhot}.} Chapter 8, mishnah 3. Hebrew text (Torat Emet; public domain via Sefaria).
\item \textbf{Mishnah, Tractate \emph{Sanhedrin}.} Chapter 6, mishnah 1. Sefaria Community Translation (CC0), checked alongside Hebrew Torat Emet text.
\item \textbf{Smith, William, ed.} \emph{A Dictionary of Greek and Roman Antiquities} (London: John Murray, 1875; LacusCurtius transcription; public domain). Entry s.v. \emph{Mantele}.
\item \textbf{Suetonius Tranquillus, Gaius.} \emph{De Vita Caesarum} (\emph{The Lives of the Caesars}), \emph{Nero} 22.2. Translated by John Carew Rolfe, Loeb Classical Library (1913--1914; LacusCurtius transcription; public domain).
\end{itemize}

\subsection{Modern Exegetical and Circulation Witnesses}

\begin{itemize}
\item \textbf{Latham, Henry.} \emph{The Risen Master: A Sequel to Pastor Pastorum}, sixth thousand (Cambridge: Deighton, Bell and Co., 1905; first published 1901; public domain), pp. 34 and 43.
\item \textbf{Westcott, Brooke Foss.} \emph{The Gospel according to St. John: The Greek Text with Introduction and Notes}, edited by Arthur Westcott, Vol. 2 (London: John Murray, 1908; public domain), p. 339 (including Arthur Westcott's editorial notes recording Bishop Westcott's marginalia and letter to H. O. Barratt).
\item \textbf{Bivin, David.} ``Christ's Linen Napkin,'' \emph{Jerusalem Perspective}, site publication date 7 October 2006 (consulted via Internet Archive capture of 7 March 2013; copyrighted web publication, cited analytically only).
\item \textbf{RecipesAndMore Mailing List Archive.} Forwarded message on ``the folded napkin,'' dated 21 May 2007 (\texttt{mail-archive.com}, message 14227; cited analytically as earliest independently archived circulation copy located).
\item \textbf{TruthOrFiction.} ``The Meaning of the Folded Napkin in Jesus' Tomb,'' inline update 28 January 2008 / byline 17 March 2015 (copyrighted web publication, cited analytically only).
\item \textbf{GotQuestions.} ``What is the significance of the folded napkin in Christ's tomb?'' (consulted July 2026; copyrighted web publication, cited analytically as a secondary transmission witness).
\end{itemize}
